% !TEX program = lualatex
% Transparencias de Fundamentos de las Matemáticas I
% Clase 15: Naturales e inducción

\documentclass[aspectratio=169,11pt]{beamer}

\usepackage{fontspec}
\usepackage[spanish,es-nodecimaldot]{babel}
\usepackage{amsmath,amssymb,mathtools}
\usepackage{booktabs}
\usepackage{csquotes}
\usepackage{xcolor}

\usetheme{Madrid}
\usecolortheme{default}
\setbeamertemplate{navigation symbols}{}
\setbeamertemplate{footline}[frame number]

\definecolor{FdMBlue}{RGB}{20,55,100}
\definecolor{FdMRed}{RGB}{130,30,30}
\definecolor{FdMGreen}{RGB}{25,100,60}

\setbeamercolor{structure}{fg=FdMBlue}
\setbeamercolor{title}{fg=white,bg=FdMBlue}
\setbeamercolor{frametitle}{fg=white,bg=FdMBlue}
\setbeamercolor{block title}{fg=white,bg=FdMBlue}
\setbeamercolor{block body}{bg=blue!3}
\setbeamercolor{alerted text}{fg=FdMRed}

\setmainfont{Latin Modern Roman}
\setsansfont{Latin Modern Sans}
\setmonofont{Latin Modern Mono}

\newcommand{\N}{\mathbb{N}}
\newcommand{\Z}{\mathbb{Z}}
\newcommand{\Q}{\mathbb{Q}}
\newcommand{\R}{\mathbb{R}}
\newcommand{\C}{\mathbb{C}}
\newcommand{\Pow}{\mathcal{P}}
\newcommand{\id}{\operatorname{id}}
\newcommand{\mcd}{\operatorname{mcd}}
\newcommand{\mcm}{\operatorname{mcm}}

\title[FdM I -- Clase 15]{Naturales e inducción}
\subtitle{Principio de inducción, inducción fuerte y buen orden}
\author{Antonio Falcó Montesinos}
\institute{Universidad CEU Cardenal Herrera}
\date{Fundamentos de las Matemáticas I}

\begin{document}

\begin{frame}
  \titlepage
\end{frame}

\begin{frame}{Objetivos de la clase}
\begin{itemize}
  \item Fijar la convención \(\mathbb{N}=\{0,1,2,\ldots\}\).
  \item Aplicar el principio de inducción.
  \item Relacionar inducción fuerte y buen orden.
\end{itemize}
\end{frame}

\begin{frame}{Mapa de la clase}
\tableofcontents
\end{frame}

\section{Naturales}

\begin{frame}{Naturales}
\begin{block}{Convención}
\begin{itemize}
  \item En este curso, \(\mathbb{N}=\{0,1,2,3,\ldots\}\).
  \item Los naturales positivos se denotan por \(\mathbb{N}^{\ast}\).
  \item La convención debe recordarse cuando el caso inicial importa.
\end{itemize}
\end{block}
\end{frame}

\begin{frame}{Naturales}
\begin{block}{Sucesor y orden}
\begin{itemize}
  \item El sucesor de \(n\) es \(n+1\).
  \item El orden \(\leq\) en \(\mathbb{N}\) es total.
  \item Todo subconjunto no vacío de \(\mathbb{N}\) tiene mínimo.
\end{itemize}
\end{block}
\end{frame}

\section{Inducción}

\begin{frame}{Inducción}
\begin{block}{Principio de inducción}
\begin{itemize}
  \item Caso base: se prueba \(P(0)\).
  \item Paso inductivo: se prueba que \(P(n)\Rightarrow P(n+1)\).
  \item Conclusión: \(P(n)\) vale para todo \(n\in\mathbb{N}\).
\end{itemize}
\end{block}
\end{frame}

\begin{frame}{Inducción}
\begin{block}{Ejemplo modelo}
\begin{itemize}
  \item Para todo \(n\in\mathbb{N}\): \(0+1+\cdots+n=\frac{n(n+1)}{2}\).
  \item El caso base es \(n=0\).
  \item El paso inductivo añade \(n+1\) a ambos lados de la fórmula conocida.
\end{itemize}
\end{block}
\end{frame}

\begin{frame}{Inducción}
\begin{block}{Inducción fuerte}
\begin{itemize}
  \item En el paso inductivo se pueden usar todos los casos anteriores.
  \item Es útil cuando \(P(n+1)\) depende de más de un caso previo.
  \item No es más verdadera: es equivalente a la inducción ordinaria en \(\mathbb{N}\).
\end{itemize}
\end{block}
\end{frame}

\section{Cierre de FdM I}

\begin{frame}{Cierre de FdM I}
\begin{block}{Síntesis}
\begin{itemize}
  \item Lenguaje preciso.
  \item Lógica y cuantificadores.
  \item Conjuntos, aplicaciones y relaciones.
  \item Estructuras algebraicas y razonamiento inductivo.
\end{itemize}
\end{block}
\end{frame}

\section{Profundización}

\begin{frame}{Profundización}
\begin{block}{Estructura de una prueba por inducción}
\begin{itemize}
  \item Nombrar la propiedad \(P(n)\).
  \item Probar el caso base correspondiente.
  \item Formular claramente la hipótesis inductiva.
  \item Demostrar el caso siguiente y cerrar con el principio de inducción.
\end{itemize}
\end{block}
\end{frame}

\begin{frame}{Profundización}
\begin{block}{Inducción fuerte: cuándo aparece}
\begin{itemize}
  \item Problemas de descomposición.
  \item Propiedades en las que el caso actual depende de varios casos anteriores.
  \item Demostraciones relacionadas con factorización o algoritmos recursivos.
\end{itemize}
\end{block}
\end{frame}

\begin{frame}{Profundización}
\begin{block}{Buen orden y mínimo contraejemplo}
\begin{itemize}
  \item Para probar una afirmación universal, se puede suponer que existe un contraejemplo.
  \item Por buen orden, habría un contraejemplo mínimo.
  \item Si se demuestra que ese mínimo no puede existir, queda probada la afirmación.
\end{itemize}
\end{block}
\end{frame}

\begin{frame}{Cierre}
\begin{itemize}
  \item La inducción es una herramienta de demostración, no una comprobación de ejemplos.
  \item El caso base depende de la convención sobre \(\mathbb{N}\).
  \item FdM I cierra con el lenguaje necesario para construir sistemas numéricos.
\end{itemize}
\end{frame}

\begin{frame}{Trabajo recomendado}
\begin{itemize}
  \item Releer las secciones correspondientes del manual.
  \item Identificar definiciones, símbolos nuevos y enunciados principales.
  \item Preparar dudas y ejemplos para el seminario asociado.
\end{itemize}
\end{frame}

\end{document}
